\section{Conclusions}\label{sec:conclusions}

In this paper, we propose WaveletMask to improve recurrent old film restoration by replacing the raw pixel-domain gate cue with a wavelet-domain gate that captures local high-frequency damage responses and global flicker responses. On the unified SRWOV re-training benchmark, WaveletMask attains the best PSNR among ten uniformly re-trained methods while adding only 898 detector parameters and remaining competitive on the perceptual metrics, and a strict paired comparison against the RTN confirms a statistically significant clip-level PSNR improvement. Additional robustness experiments across three random seeds and artifact-stratified analyses further show that the wavelet-domain gate is particularly effective for scratch- and flicker-related degradations. Beyond full-reference evaluation, the method also shows practical value on real archival clips: WaveletMask is the only method placing within the top two on all three no-reference metrics (narrowly so on CLIPIQA+), and the qualitative comparisons show cleaner scratch and flicker removal with better-preserved film texture and temporal stability. Taken together, these findings indicate that WaveletMask provides a lightweight and effective degradation-aware gating strategy for old film restoration. Because the wavelet-domain gate adds only 898 detector parameters on top of an existing recurrent backbone, it can be integrated as a drop-in gating cue into current restoration pipelines, making it practical for large-scale archival-film and legacy-video remastering where per-clip overhead matters.

These gains nonetheless come with clear boundary conditions: persistent image-plane-fixed scratches, abrupt temporal discontinuities such as scene cuts, and spatially extensive film losses remain challenging for the current restoration-only formulation (Section~\ref{subsec:limitations}). Future work can further improve persistent scratch handling and replace the current fusion design with learned fusion under explicit supervision. It is also important to evaluate the method on more diverse real-world archival datasets to strengthen cross-degradation generalization and robustness in challenging restoration scenarios.
